\title{Direct Preference Optimization: Your Language Model is Secretly a Reward Model}

\begin{abstract}
Large-scale unsupervised language models (LMs) acquire extensive general knowledge and exhibit some capacity for reasoning; however, effectively steering their outputs remains challenging due to the entirely unsupervised approach used during pretraining. Prevailing strategies to increase steerability typically involve gathering human annotations comparing the quality of different outputs and then fine-tuning the LM to reflect these judgments, most frequently through reinforcement learning from human feedback (RLHF). Nevertheless, RLHF is a multifaceted and often unstable process that first trains a reward model mirroring human preferences and subsequently optimizes the original LM using reinforcement learning to maximize this surrogate reward, while constraining divergence from the pretrained model. This work presents a novel formulation of the reward model within the RLHF framework that permits derivation of the associated optimal policy in a closed-form expression, reducing the RLHF process to minimization of a straightforward classification objective. Our method, termed \textit{Direct Preference Optimization} (DPO), achieves strong empirical results, exhibiting robustness and efficiency in computation, as it obviates the requirement for language model sampling during fine-tuning and reduces the burden of hyperparameter selection. Empirical studies demonstrate that DPO matches or outperforms existing techniques in aligning LMs with human preferences. Specifically, DPO surpasses PPO-based RLHF in regulating the sentiment of generated text and achieves comparable or better response quality in summarization and single-turn dialogue tasks, all while being notably easier to implement and train.
\end{abstract}

\section{Introduction}
Large-scale unsupervised language models (LMs) trained on massive corpora have demonstrated remarkable emergent abilities~\citep{chowdhery2022palm, brown2020language, touvron2023llama, bubeck2023sparks}. Nevertheless, the data fueling their training originates from a diverse array of human-authored sources, reflecting an assortment of intentions, skill levels, and objectives. Many of these underlying human motivations and competencies are not always desirable for imitation; for instance, although it is beneficial for an AI coding assistant to \textit{recognize} prevalent programming mistakes to enable correction, we would still prefer that the model’s code generation predominantly emulate the superior—if comparatively uncommon—coding standards exemplified within its training set. Analogously, a language model should ideally be \textit{informed} of a widespread misconception affecting 50\% of individuals, yet we do not wish the model to affirm this misconception as factual in 50\% of related interactions. Thus, the capability to delineate a model’s \emph{targeted responses and conduct} from its broad \textit{knowledge base and skill set} is fundamental for developing AI that is reliable, high-performing, and governable \citep{ouyang2022training}. While prevalent strategies employ reinforcement learning (RL) to align LMs with human value judgments, our work demonstrates that the objective currently used in RL-based approaches can, in fact, be exactly realized using a simple binary cross-entropy loss, thereby significantly streamlining the preference learning process.

In broad terms, contemporary approaches to preference alignment embed desirable patterns of behavior into a language model through curated datasets encoding human preference signals concerning safety and helpfulness. This phase of preference optimization is subsequent to initial extensive unsupervised pre-training on large-scale textual datasets. Although one straightforward method for preference modeling involves supervised fine-tuning using exemplary human-authored responses, the dominant and most effective methodology remains reinforcement learning from human (or AI) feedback (RLHF/RLAIF; \citep{christiano2017deep,bai2022constitutional}). In RLHF frameworks, a reward function is learned from a dataset of comparative human judgments, after which RL is applied to adjust the model's policy to elicit high-reward responses, all while preventing substantial divergence from the pre-trained policy. Despite RLHF’s success in producing models with state-of-the-art dialogue and programming abilities, the underlying pipeline is notably more intricate than standard supervised training, as it necessitates the training of several models and iterative sampling from the policy during optimization, resulting in elevated computational demands.

This work demonstrates a method for optimizing language models in line with human preferences, circumventing the need for explicit reward modeling or reinforcement learning frameworks. We introduce \textit{Direct Preference Optimization (DPO)}, an approach that implicitly targets the same objective as standard RLHF techniques—namely, reward maximization under a KL-divergence constraint—while offering a simpler implementation and training procedure. Conceptually, the DPO update boosts the log-probability ratio of preferred versus non-preferred outputs, yet it utilizes a dynamic, example-specific importance weighting scheme. This mechanism helps avoid the collapse issues observed with straightforward probability ratio objectives. Similar to prior approaches, DPO employs a theoretical preference framework (such as the Bradley-Terry model; \cite{bradley1952rankanalysis}) to evaluate the agreement between reward functions and observed preference data. In contrast to other algorithms, which define a preference-based loss to train a reward model and subsequently a policy to optimize this learned reward, DPO employs a variable transformation to express the preference loss directly as a function of the policy. Using human-annotated preference data over model outputs, DPO is able to optimize the policy directly via a standard binary cross-entropy loss, yielding the optimal policy for an implicitly derived reward function reflecting the preference data.

The central contribution of this work is Direct Preference Optimization (DPO), a straightforward, reinforcement learning-free technique for fine-tuning language models using human preferences. Empirical results indicate that DPO matches or surpasses the performance of prevalent methods, including RLHF implementations based on PPO, across various preference-learning tasks such as sentiment control, summarization, and conversation, even when scaling to language models with up to 6B parameters.

\section{Related Work}

Self-supervised language models at larger scales have demonstrated the ability to handle some tasks without explicit training examples (zero-shot) \citep{radford2019language}, or with a limited number of examples (few-shot prompts) \citep{gpt3,megatron,chowdhery2022palm}. Nonetheless, substantial gains in both performance on downstream applications and in adhering to user intent can be achieved by further training on datasets containing explicit instructions and human-generated responses \citep{mishra-etal-2022-cross,sanh2022multitask,chung2022scaling,thoppilan2022lamda}. This technique, known as `instruction-tuning', equips large language models (LLMs) with the capacity to generalize beyond the set of instructions seen during fine-tuning, thereby enhancing their overall practical effectiveness \citep{chung2022scaling}. Although instruction tuning has proven effective, it is frequently less labor-intensive to gather comparative human feedback on response quality than to amass expert-curated examples. Consequently, recent research has focused on refining LLMs using datasets that reflect human preference rankings, leading to improved capabilities in tasks such as translation \citep{kreutzer-etal-2018-reliability}, summarization \citep{stiennon2022learning,ziegler2020finetuning}, story generation \citep{ziegler2020finetuning}, and compliance with instructions \citep{ouyang2022training,ramamurthy2023is}. This approach first involves training a neural reward model to fit the collected preference data, often via a preference model like the Bradley-Terry model \citep{bradley1952rankanalysis}, and subsequently applying reinforcement learning algorithms—such as REINFORCE \citep{williams1992reinforce}, proximal policy optimization (PPO; \cite{schulman2017proximal}), or their derivatives \citep{ramamurthy2023is}—to maximize the defined reward in the language model. Related work has utilized instruction-following LLMs, already aligned using human feedback, to synthesize additional preference-labeled data targeting specific characteristics like safety or non-toxicity \citep{bai2022constitutional}, relying primarily on minimal human supervision in the form of rubric-based guidance for LLM annotations. These developments illustrate a merging of two research directions: reinforcement learning-based language model training for diverse goals~\citep{Ranzato2015SequenceLT,paulus2018a,wu2018learning}, and methodologies focused on learning from human preference information \citep{christiano2017deep,kupcsik2018learning}. In spite of the promise shown by preference-based approaches, applying reinforcement learning at scale for LLM fine-tuning presents significant practical difficulties; this work introduces a theoretically-grounded alternative for optimizing relative preferences that circumvents the need for reinforcement learning.

Beyond language-related applications, the study of policy learning from preferences has been explored within both bandit frameworks and reinforcement learning, giving rise to a range of methodological approaches. In contextual bandit problems, the paradigm of learning from preferences or action rankings, rather than explicit rewards, is referred to as the contextual dueling bandit (CDB) problem \citep{yue2012karmed, dudik2015contextual}. In scenarios lacking absolute reward feedback, the analysis of CDBs replaces the concept of an optimal policy with that of a \textit{von Neumann winner}—that is, a policy whose expected probability of outperforming any competing policy meets or exceeds 50\% \citep{dudik2015contextual}. Nevertheless, a distinguishing aspect of the CDB setting is that preference feedback is provided online, whereas typical human preference learning involves static, offline datasets consisting of preference-annotated action pairs \citep{yan2022human}. Analogously, \textit{preference-based RL} (PbRL) addresses learning from binary preference signals generated by an unknown 'scoring' function in lieu of explicit rewards \citep{BusaFekete2014, ruiz2023dueling}. Several PbRL algorithms have been developed—including those capable of utilizing off-policy preference data—but most require a two-step procedure: first estimating the latent scoring (reward) function, and then optimizing the policy relative to this surrogate objective \citep{jain2013learning, BusaFekete2014, christiano2017deep, sadigh2017active, kupcsik2018learning}. By contrast, we introduce a unified approach in which a policy is trained in a single stage to directly satisfy preferences.

\section{Preliminaries}\label{section:prelims}

We provide an overview of the RLHF process as described by \citeauthor{ziegler2020finetuning} and extended in works such as \citep{stiennon2022learning, bai2022training, ouyang2022training}. This pipeline typically comprises three sequential stages: (1) supervised fine-tuning (SFT); (2) collecting preference data and training a reward model; and (3) reinforcement learning-based optimization.

\textbf{SFT}: The RLHF framework generally starts by adapting a pre-trained language model (LM) using supervised learning. The model is trained on curated datasets relevant to the downstream application (such as dialogue systems or summarization), resulting in an adapted model denoted $\pi^\text{SFT}$.

\textbf{Reward Modelling Phase}: In the next step, the SFT model receives prompts $x$ and generates pairs of candidate responses $(y_1, y_2)\sim \pi^\text{SFT}(y \mid x)$. Human evaluators then review these pairs, indicating their preferred answer; a preference annotation is represented as $y_w\succ y_l \mid x$, where $y_w$ is favored over $y_l$ for prompt $x$. These choices are presumed to arise from an underlying but unobservable reward model $r^*(y, x)$. Several statistical models for inferring preferences exist; one commonly used option is the Bradley-Terry (BT) model \cite{bradley1952rankanalysis}—although the framework also accommodates more general models, such as the Plackett-Luce model \citep{plackett1975analysis, luce2012individual}, particularly when handling lists with multiple ranked outputs. According to the BT formulation, the human preference distribution $p^*$ can be described as:

Given a fixed dataset of preference pairs $\mathcal{D}=\bigl\{x^{(i)}, y_w^{(i)}, y_l^{(i)}\bigr\}_{i=1}^N$ sampled from $p^*$, we can represent the reward with a parameterized model $r_{\phi}(x, y)$, and estimate its parameters by maximizing the likelihood. Treating the task as binary classification leads to a negative log-likelihood objective of

\begin{equation}\label{eq:reward_model}
    \mathcal{L}_R(r_{\phi}, \mathcal{D}) = -\mathbb{E}_{(x, y_w, y_l)\sim \mathcal{D}}\bigl[\log \sigma(r_{\phi}(x, y_w)- r_{\phi}(x, y_l))\bigr],
\end{equation}

where $\sigma$ is the sigmoid activation. Within the language modeling framework, the function $r_{\phi}(x, y)$ commonly inherits weights from the SFT policy $\pi^\text{SFT}(y \mid x)$, augmented with a linear projection attached to the terminal transformer layer to yield scalar rewards~\cite{ziegler2020finetuning}. To further decrease the variance in reward predictions, existing literature normalizes $r_\phi$ so that $\mathbb{E}_{x,y\sim \mathcal{D}}\left[r_\phi(x, y)\right] = 0$ holds for all $x$.

\textbf{RL Fine-Tuning Phase}: In this stage, the model trained to approximate the reward is used to evaluate outputs from the language model, and learning is then driven by these scores. As detailed by earlier works~\citep{jaques2017sequence, jaques2020human}, the fine-tuning objective can be expressed as

\begin{equation}\label{eq:RL}
\max_{\pi_{\theta}}  \mathbb{E}_{x\sim \mathcal{D}, y\sim \pi_{\theta}(y \mid x)}\bigl[r_{\phi}(x, y)\bigr] - \beta\mathbb{D}_{\textrm{KL}}\bigl[\pi_{\theta}(y\mid x)\mid \mid \pi_\text{ref}(y\mid x)\bigr],
\end{equation}

with $\beta$ acting as a scaling factor restricting divergence from a reference distribution, typically taken as the initial SFT model $\pi^\text{SFT}$. Generally, $\pi_{\theta}$ starts from $\pi^\text{SFT}$ as well. The inclusion of this penalty term serves to regularize the policy update, ensuring it remains close to the support of data where the reward estimates are reliable, and helping to preserve diversity and mitigate collapse into a narrow set of outputs. Since language generation operates in a discrete space, direct gradient-based optimization is infeasible, necessitating reinforcement learning algorithms. Current practice~\citep{ziegler2020finetuning, stiennon2022learning, bai2022training, ouyang2022training} reformulates the objective in terms of the reward function as $r(x, y) = r_{\phi}(x, y) -\beta (\log \pi_{\theta}(y\mid x) - \log \pi_\text{ref}(y\mid x))$, with optimization performed by PPO~\cite{schulman2017proximal}.

\section{Direct Preference Optimization}\label{sec:DPO}

Driven by the difficulties encountered when applying standard RL techniques to large-scale language model fine-tuning, we propose a streamlined method for policy learning that directly utilizes human preference data. Rather than following traditional RLHF workflows—wherein a reward model is trained and then optimized via reinforcement learning—our technique employs a special parameterization for the reward model, facilitating recovery of the optimal policy analytically, without necessitating iterative RL procedures. As we elaborate below, our central insight is to exploit a closed-form relation linking reward models to their corresponding optimal policies. This relationship lets us transform objectives expressed over reward functions into equivalent losses directly defined on the policy itself. By leveraging this change of variables, we circumvent the need for a separately trained reward predictor while maintaining consistency with standard preference modeling approaches such as the Bradley-Terry framework. In this formulation, the policy network inherently encapsulates both the language modeling and the (implicit) reward estimation.

\textbf{Deriving the DPO objective.} We begin with the RL objective described in Eq.~\ref{eq:RL}, employing a general reward function $r$, as established in previous literature. Building on earlier studies~\citep{peters2007reinforcement, peng2019advantage, korbak2022reinforcement, go2023aligning}, it can be readily verified that the solution optimizing the KL-regularized reward maximization problem in Eq.~\ref{eq:RL} adopts the following structure:

\begin{equation}\label{eq:op_policy}
    \pi_r(y\mid x) = \frac{1}{Z(x)}\pi_\text{ref}(y\mid x)\exp\left(\frac{1}{\beta}r(x, y)\right),
\end{equation}

Here, $Z(x) =\sum_{y}\pi_\text{ref}(y\mid x)\exp\left(\frac{1}{\beta}r(x, y)\right)$ serves as the partition function. The details of the derivation appear in Appendix \ref{app:derivation1}. While employing the MLE estimate $r_{\phi}$ for the underlying true reward $r^*$ is feasible, the practical estimation of $Z(x)$ is computationally intensive~\citep{korbak2022reinforcement, go2023aligning}, presenting obstacles for direct application. Nonetheless, Eq.~\ref{eq:op_policy} can be rearranged to express the reward in terms of the associated optimal policy $\pi_r$, the reference policy $\pi_\text{ref}$, and the partition function $Z(\cdot)$. Specifically, by applying the logarithm to both sides of Eq.~\ref{eq:op_policy} and carrying out algebraic simplification, we obtain:

\begin{equation}\label{eq:main_eq}
    r(x,y) =\beta \log \frac{\pi_r(y\mid x)}{\pi_\text{ref}(y\mid x)} + \beta \log Z(x).
\end{equation}

This change of variables can be applied to both the ground-truth reward $r^*$ and its associated optimal policy $\pi^*$. Notably, the Bradley-Terry model relies solely on the relative difference between the rewards of two candidate outputs, that is, ${p^*(y_1 \succ y_2 \mid x) = \sigma(r^*(x, y_1) - r^*(x, y_2))}$. Plugging the substitution from Eq.~\ref{eq:main_eq} for $r^*(x,y)$ into the preference framework in Eq.~\ref{eq:bradley-terry}, the partition functions $Z(x)$ drop out, allowing the human preference likelihood to be expressed strictly through the optimal policy $\pi^*$ and the reference policy $\pi_\text{ref}$. Accordingly, the Bradley-Terry-based RLHF optimal policy $\pi^*$ fulfills:

\begin{equation}\label{eq:objective}
    p^*(y_1\succ y_2 \mid x)=\frac{1}{1 + \exp\left(\beta \log \frac{\pi^*(y_2\mid x)}{\pi_\text{ref}(y_2\mid x)} - \beta \log \frac{\pi^*(y_1\mid x)}{\pi_\text{ref}(y_1\mid x)}\right)}
\end{equation}

The complete derivation is presented in Appendix~\ref{app:derivation2}. Although the form in Eq.~\ref{eq:objective} is based on the Bradley-Terry model, similar formulations arise for broader Plackett-Luce models~\citep{plackett1975analysis, luce2012individual}, as explained in Appendix~\ref{app:plackett_luce_models}.

Having established an expression for the human preference probabilities using only the optimal policy, we can proceed to specify a maximum likelihood learning objective for a parameterized policy $\pi_\theta$. Mirroring the methodology of reward modeling (as in Eq.~\ref{eq:reward_model}), the corresponding policy objective is given by:

\begin{equation}\label{eq:optimum_model}
    \mathcal{L}_\text{DPO}(\pi_{\theta}; \pi_\text{ref}) = -\mathbb{E}_{(x, y_w, y_l)\sim \mathcal{D}}\left[\log \sigma \left(\beta \log \frac{\pi_{\theta}(y_w\mid x)}{\pi_\text{ref}(y_w\mid x)} - \beta \

In this approach, we model an implicit reward function through an alternative parameterization, for which the optimal policy corresponds directly to $\pi_\theta$. Additionally, because our methodology effectively reparameterizes the Bradley-Terry model, it inherits theoretical guarantees such as consistency, provided that the preference data distribution satisfies appropriate conditions \cite{bong2022generalized}. A more detailed exploration of these theoretical aspects in the context of DPO and related frameworks is provided in Section~\ref{sec:theory}.

\textbf{How does the DPO update function?} To gain mechanistic insight into DPO, one can examine the gradient of its loss function, $\mathcal{L}_\text{DPO}$. The parameter gradient can be formulated as:

\begin{multline*}\label{eq:gradient}
    \nabla_\theta \mathcal{L}_\text{DPO}(\pi_\theta;\pi_\text{ref}) = \\ -\beta\mathbb{E}_{(x, y_w, y_l) \sim \mathcal{D}} \bigg[\underbrace{\sigma(\hat{r}_\theta(x, y_l) - \hat{r}_\theta (x, y_w))}_\text{greater contribution when the reward is inaccurately estimated}\bigg[\underbrace{\nabla_\theta\log \pi(y_w \mid x)}_\text{promote $y_w$} - \underbrace{\nabla_\theta\log\pi(y_l \mid x)}_\text{demote $y_l$}\bigg]\bigg],
\end{multline*}

with $\hat{r}_\theta(x, y) = \beta \log \frac{\pi_\theta(y \mid x)}{\pi_\text{ref}(y \mid x)}$ denoting the implicit reward as derived from the language model $\pi_\theta$ relative to the reference model $\pi_\text{ref}$ (see also Section~\ref{sec:theory}). In practical terms, the gradient of $\mathcal{L}_\text{DPO}$ encourages higher likelihood for preferred completions $y_w$, while simultaneously reducing the likelihood assigned to non-preferred completions $y_l$. Crucially, each training example is scaled according to the extent to which the implicit reward model $\hat{r}_\theta$ overvalues the non-preferred outputs, modulated by $\beta$, reflecting the severity of reward misordering as well as the influence of the KL regularization. Empirically, this specific weighting is critical—an ablation that omits the weighting factor can result in pathological degeneration of the language model (see Appendix Table~\ref{tab:unlikelihood_generations}).

\textbf{DPO framework.} The typical DPO workflow includes the following steps: 1) For each prompt $x$, generate candidate completions $y_1, y_2 \sim \pi_\text{ref}(\cdot \mid x)$, assign human preference labels, and construct an offline preference dataset $\mathcal{D} = \{x^{(i)}, y_w^{(i)}, y_l^{(i)}\}_{i=1}^N$; 2) Update the language model $\pi_\theta$ to minimize $\mathcal{L}_\text{DPO}$ using the fixed $\pi_\text{ref}$, dataset $\mathcal{D}$, and a chosen $\beta$. 
In real-world applications, it is often desirable to leverage existing publicly available preference datasets instead of collecting new samples and annotations. Since these datasets are generated using $\pi^\text{SFT}$, we set $\pi_\text{ref} = \pi^\text{SFT}$ whenever this reference model is accessible. When $\pi^\text{SFT}$ is not accessible, we instead construct $\pi_\text{ref}$ by maximizing the likelihood over the preferred completions $(x, y_w)$ from the data, i.e., $\pi_\text{ref} = \argmax_{\pi}\mathbb{E}_{x, y_w \sim \mathcal{D}}\left[\log \pi(y_w \mid x)\right]$. This approach helps address the discrepancy between the intractable true reference distribution and the practical $\pi_\text{ref}$ used in DPO. Additional implementation specifics and information about hyperparameters are provided in Appendix~\ref{app:implementation}.

\section{Theoretical Analysis of DPO} In this section, we provide a deeper theoretical interpretation of the DPO approach, present supporting analyses, and connect DPO’s strengths to weaknesses seen in actor-critic methods commonly adopted for RLHF (e.g., PPO~\cite{schulman2017proximal}).

\label{sec:theory}

\subsection{Your Language Model Is Secretly a Reward Model} The DPO approach enables direct policy optimization via a single maximum likelihood objective, circumventing the need to first train a separate reward model and then run RL for policy learning. Importantly, the optimization criterion in Eq.~\ref{eq:main_eq} aligns with a Bradley-Terry model, where the reward takes the form $r^*(x, y) = \beta \log\frac{\pi^*_\theta(y \mid x)}{\pi_\text{ref}(y \mid x)}$. Optimizing the parameterized model $\pi_\theta$ thus corresponds—under a change of variables—to the same solution obtained by explicit reward model training as described in Eq.~\ref{eq:reward_model}. In what follows, we formalize the theoretical underpinnings of this reparameterization, demonstrate that it does not restrict the expressivity of learned reward models, and prove that the method permits exact recovery of the optimal policy. We start by introducing an equivalence relation over reward functions.

\begin{definition}
Two reward functions $r(x, y)$ and $r'(x, y)$ are said to be equivalent if and only if ${r(x, y)-r'(x, y) = f(x)}$ for some function $f$.
\end{definition}

It can be readily verified that this defines an equivalence relation, segmenting the space of reward functions into equivalence classes. We next state two important lemmas:

\begin{lemma}\label{lemma:same_prefrence}
Within the Plackett-Luce model, and specifically within the Bradley-Terry framework, any pair of reward functions belonging to the same equivalence class will produce identical preference distributions.
\end{lemma}

\begin{lemma}\label{lemma:same_policy}
    If two reward functions lie in the same equivalence class, then the induced optimal policies in the constrained RL setting will coincide.
\end{lemma}

The arguments for these statements are elementary, so we include the details in Appendix \ref{app:lemma1}. The first lemma highlights the familiar ambiguity present in the Plackett-Luce family of preference models \cite{plackett1975analysis}, often referred to as under-specification. As a consequence, identifiability constraints must typically be applied to ensure well-posedness of MLE estimators such as those in Eq. \ref{eq:reward_model} \cite{bong2022generalized}. The second lemma asserts that all members within a given equivalence class of reward functions yield the same optimal policy; thus, our primary aim reduces to finding any representative of the optimal reward class. Our main result is presented below and proven in Appendix~\ref{app:thm1}:

\begin{theorem}\label{thm:main}
    Assuming mild conditions, every reward equivalence class that is compatible with the Plackett-Luce (or, in particular, the Bradley-Terry) model admits a parameterization of the form ${r(x, y) = \beta \log \frac{\pi(y\mid x)}{\pi_\text{ref}(y\mid x)}}$ for some model $\pi(y\mid x)$ and a fixed reference model $\pi_\text{ref}(y \mid x)$.
\end{theorem}

\begin{sproof}
    Take any reward function $r(x, y)$, which determines an associated optimal model $\pi_r(y \mid x)$ as characterized by Eq. \ref{eq:op_policy}. Our goal is to demonstrate that a reward function from the equivalence class containing $r$ can be constructed via the aforementioned reparameterization. Define the following projection:
\begin{equation}
    f(r; \pi_\text{ref}, \beta)(x, y) = r(x, y) - \beta\log\sum_{y}\pi_\text{ref}(y\mid x)\exp\left(\frac{1}{\beta}r(x, y)\right)
\end{equation}
This operator $f$ normalizes the reward by subtracting the log of the partition function of $\pi_r$; notably, this normalization depends solely on $x$. Thus, $f(r; \pi_\text{ref}, \beta)(x, y)$ remains within the equivalence class of $r(x, y)$. Substituting $r$ using the right side of Eq.~\ref{eq:main_eq} (valid for all reward functions), we obtain $f(r; \pi_\text{ref}, \beta)(x, y) = \beta \log \frac{\pi_r(y\mid x)}{\pi_\text{ref}(y\mid x)}$. This shows that the projection $f$ yields a reward function of the desired parametric form, thereby establishing that the reparameterization retains full generality for our purposes.
\end{sproof}

Theorem~\ref{thm:main} can be reinterpreted as identifying the precise reward function, within each equivalence class, chosen by the DPO reparameterization. Specifically, it selects the reward function that fulfills the condition:

\begin{equation}\label{eq:lag_p}
     \sum_{y}\underbrace{\pi_\text{ref}(y\mid x)\exp\left(\frac{1}{\beta}r(x, y)\right)}_{=\pi(y\mid x)\text{, using Thm.~\ref{thm:main} reparam.}} = 1,
\end{equation}

meaning that $\pi(y\mid x)$ constitutes a proper probability distribution (with non-negative probabilities that sum to one). Moreover, by referring to Eq.~\ref{eq:op_policy}, it becomes evident that Eq.~\ref{eq:lag_p} expresses the partition function for the optimal policy defined by the reward $r(x, y)$. The principal insight underpinning the DPO algorithm is the observation that one can introduce specific constraints into the otherwise under-determined Plackett-Luce family (including the Bradley-Terry model) of preference frameworks. In doing so, we retain the expressive power of the original reward model family while ensuring that the optimal policy from Eq.~\ref{eq:op_policy} becomes analytically solvable for any prompt $x$.

\subsection{Instability of Actor-Critic Algorithms}
Our framework also provides a lens to analyze the instability issues common in actor-critic methods employed in RLHF, such as PPO. We remain within the RLHF process, specifically concentrating on the RL fine-tuning phase as delineated in Section \ref{section:prelims}. By connecting this to the control as inference framework \cite{levine2018reinforcement}, for the constrained RL formulation in Eq.~\ref{eq:RL}, we posit a parameterized distribution $\pi_{\theta}(y\mid x)$, seeking to minimize $\mathbb{D}_{\text{KL}}[\pi_{\theta}(y|x) \mid \mid \pi^*(y\mid x)]$, where $\pi^*$ is the optimal policy derived from Eq.~\ref{eq:optimum_model} and the associated reward function $r_{\phi}(y, x)$. After some algebraic manipulation, this process yields the following objective:

\begin{equation}\label{eq:AC}
    \max_{\pi_{\theta}}\mathbb{E}_{\pi_{\theta}(y\mid x)}\bigg[\underbrace{r_{\phi}(x, y) -\beta\log\sum_{y}\pi_\text{ref}(y\mid x)\exp\left(\frac{1}{\beta}r_{\phi}(x, y)\right)}_{f(r_{\phi}, \pi_\text{ref}, \beta)} - \underbrace{\beta\log\frac{\pi_{\theta}(y\mid x)}{\pi_\text{ref}(y\mid x)}}_{\text{KL}}\bigg]
\end{equation}

The objective discussed here aligns with those targeted in earlier studies 
\citep{ziegler2020finetuning, stiennon2022learning, bai2022training, ouyang2022training}, where optimization is performed with a DPO-equivalent reward specific to the $r_{\phi}$ reward class. In this framework, the normalization component present in $f(r_{\phi}, \pi_\text{ref}, \beta)$ can be viewed as the soft value function corresponding to the reference policy $\pi_\text{ref}$. Although this normalization does not influence the location of the optimum, omitting it may result in high variance in the policy gradient, thereby destabilizing the learning process. It is possible to address the normalization term through the use of a parameterized value function, yet this approach often encounters optimization challenges. Previous research has also employed reward normalization using human completions as baselines, which serves as a Monte Carlo estimate of the normalization term from a single sample. Distinctly, the DPO reparameterization constructs a reward function that eliminates the requirement for any baseline.

\section{Experiments}
Here, we assess DPO’s empirical effectiveness at training policies solely from preference data. To start, we consider a carefully controlled text generation environment to examine how effectively DPO balances maximizing the reward and minimizing KL-divergence to the reference policy relative to other preference learning approaches, such as PPO. Subsequently, we test DPO on larger-scale models and more complex RLHF tasks, including summarization and conversational tasks. Our results indicate that DPO frequently matches or exceeds the performance of strong baselines—including RLHF with PPO and strategies such as taking the best among $N$ samples according to a trained reward model—requiring little to no hyperparameter tuning. Prior to detailing these findings, we outline the experimental methodology; supplementary specifics can be found in Appendix~\ref{app:exp_details}.

\textbf{Tasks.} The scope of our experiments covers three open-ended text generation tasks. Across all tasks, algorithms are tasked with learning a policy from a dataset of preferences $\mathcal{D}=\bigl\{x^{(i)}, y_w^{(i)}, y_l^{(i)}\bigr\}_{i=1}^N$. In the \textbf{controlled sentiment generation} setting, the prompt $x$ is an excerpt from a movie review taken from the IMDb dataset \cite{maas-EtAl:2011:ACL-HLT2011}. The objective is for the policy to generate a continuation $y$ that expresses positive sentiment. For a controlled experimental setup, preference pairs between generated outputs are synthesized using an existing sentiment classifier, such that $p(\text{positive}\mid x,y_w)>p(\text{positive}\mid x,y_l)$. Supervised fine-tuning (SFT) involves training GPT-2-large to convergence using positive reviews from the IMDb training set, as described in more depth in App~\ref{app:sentiment_details}. For the \textbf{summarization} task, $x$ corresponds to a Reddit forum post, and the aim is for the policy to generate a summary $y$ capturing the post’s central themes. Consistent with established methodology, the Reddit TL;DR summarization dataset \citep{volske-etal-2017-tl} is used in conjunction with human preference annotations collected by \citeauthor{stiennon2022learning}. Here, the SFT baseline is a model adapted to human-generated summaries of forum posts\footnote{\url{https://huggingface.co/CarperAI/openai_summarize_tldr_sft}}, trained with the TRLX \citep{leandro_von_werra_2023_7790115} toolkit for RLHF. Notably, the human-labeled preference dataset is derived from outputs of a related, but separately trained, SFT model. The third task, \textbf{single-turn dialogue}, sets $x$ as a user message, ranging from scientific queries to personal advice. The model must provide a useful and engaging response $y$; this is evaluated on the Anthropic Helpful and Harmless dialogue dataset \citep{bai2022training}, featuring 170,000 human-assistant dialogues. Each conversation concludes with two model-generated responses and a human judgment of preference. In this domain, no specialized SFT model exists, so we construct one by fine-tuning a publicly available language model on the responses labeled as preferred.

\textbf{Evaluation.} We adopt two primary strategies for evaluating our models. To measure the capability of each method in maximizing the constrained reward objective, we conduct experiments in a controlled sentiment generation environment. Here, we assess each algorithm by plotting the trade-off between attained reward and KL-divergence relative to the reference policy; this analysis is enabled by access to a ground-truth reward function, operationalized as a sentiment classifier. Outside this synthetic setting, real-world evaluation is complicated by the lack of an explicit reward function. Accordingly, we turn to \textit{win rate} metrics against a baseline policy, using GPT-4 as a stand-in for human raters to evaluate summary quality and response helpfulness in summarization and single-turn dialogue tasks, respectively. For the summarization domain, baseline comparisons are drawn from test set reference summaries, while for dialogue, they are based on the preferred responses found in the dataset. Although prior work indicates that large language models may serve as superior automatic evaluators compared to standard metrics \citep{Chen2023ExploringTU}, we further validate our reliance on GPT-4 via a human annotation study described in Sec.~\ref{sec:human-judgments}. Results show a high degree of concordance between GPT-4 and human judgments, with humans often matching or surpassing inter-annotator agreement when compared to GPT-4.

\textbf{Methods.} Alongside DPO, our analysis encompasses a range of prevailing methodologies for aligning language model outputs with human preferences. For the summarization task, we investigate zero-shot prompting using \textbf{GPT-J} \citep{gpt-j}; for dialogue, we apply 2-shot prompting to \textbf{Pythia-2.8B} \citep{biderman2023pythia}. Our evaluation also includes the \textbf{SFT} baseline as well as \textbf{Preferred-FT}, a model obtained through supervised fine-tuning on the preferred completions $y_w$—either sourced from the SFT model in sentiment and summarization settings or from a standard language model in dialogue. We further examine a pseudo-supervised strategy: \textbf{Unlikelihood}~\citep{welleck2019neural}, which involves maximizing the likelihood of $y_w$ while explicitly discouraging the likelihood of $y_l$, using an optional weighting factor $\alpha\in[0,1]$ for the unlikelihood component. Our comparison set also includes \textbf{PPO} \citep{schulman2017proximal}, trained with rewards learned from preference datasets, and an oracle variant, \textbf{PPO-GT}, that directly leverages the known ground-truth reward in the controlled sentiment experiments. For PPO-GT, we use both a standard implementation \cite{leandro_von_werra_2023_7790115} and an adjusted version with reward normalization and enhanced hyperparameter tuning—modifications that are likewise applied to PPO when using learned rewards. To round out our baselines, we employ the \textbf{Best of $N$} approach, in which $N$ outputs are sampled from the SFT (or Preferred-FT for dialogue) model and the candidate with the highest predicted reward, according to the learned reward model, is selected. Although this method tends to perform strongly, its high computational demands render it less viable for large $N$ due to the need to generate multiple outputs per input during evaluation.

\subsection{How well can DPO optimize the RLHF objective?}

RLHF methods typically aim to maximize reward subject to a KL divergence penalty, which discourages the policy from drifting too far from a reference policy. As such, evaluations of algorithmic performance should jointly consider both achieved reward and KL divergence: a marginal reward gain accompanied by a substantial increase in KL may be suboptimal. Figure~\ref{fig:frontier-tldr-main} presents the trade-off between reward and KL across different algorithms within the sentiment classification context. For each method, we conduct several independent training experiments, varying the policy conservativeness via key hyperparameters (target KL $\in\{3,6,9,12\}$ for PPO, $\beta \in \{0.05,0.1,1,5\}$ for DPO, $\alpha\in\{0.05,0.1,0.5,1\}$ for unlikelihood, and using distinct random seeds for preferred-FT), amounting to 22 experimental runs in total. After every 100 training steps, up to convergence, we assess each trained policy on a test suite, computing both the average reward according to the actual reward model and the mean sequence-level KL\footnote{This is calculated as the sum of KL divergences across all time steps.} with respect to the reference policy $\text{KL}\left(\pi\mid \mid \pi_\text{ref}\right)$. Our analysis demonstrates that DPO achieves a superior reward-KL Pareto frontier compared to the other algorithms, attaining the highest rewards while maintaining relatively low KL divergence. This outcome is particularly striking for several reasons. To begin with, both DPO and PPO optimize an identical objective, yet DPO attains this with substantially greater sample efficiency—its reward/KL tradeoff consistently outperforms that of PPO. Furthermore, DPO outperforms PPO even when PPO is provided with access to true reward signals (PPO-GT).

\subsection{Can DPO scale to real preference datasets?}
\label{sec:dpo-real-datasets}
We proceed to assess the capability of DPO to fine-tune language models on tasks such as summarization and single-turn dialogue, utilizing real-world preference data. In the case of summarization, automatic metrics like ROUGE have been shown to correlate poorly with human judgments~\citep{stiennon2022learning}. Previous research suggests that optimizing language models with PPO according to human feedback can yield more preferred summaries. In our experiments, we compare various approaches by generating completions on the test portion of the TL;DR summarization dataset, then calculating each method's average win rate against the dataset's reference summaries. For all models, outputs are sampled at temperature settings between 0.0 and 1.0, with resulting win rates illustrated in Figure~\ref{fig:frontier-tldr-main} (right). All methods—including DPO, PPO, and Preferred-FT—start from the same GPT-J SFT model\footnote{\url{https://huggingface.co/CarperAI/openai_summarize_tldr_sft}}. Our findings indicate that DPO achieves a win rate close to 61\% at a sampling temperature of 0.0, surpassing PPO, which attains about 57\% under its best-performing temperature. Furthermore, DPO reaches a higher peak win rate compared to the best of $N$ baseline method. It is worth mentioning that we did not extensively optimize the $\beta$ hyperparameter for DPO, so these outcomes might not fully represent the method’s capabilities. Additionally, DPO shows greater resilience to changes in sampling temperature relative to PPO, which suffers performance drops at higher temperatures—sometimes equating to the base GPT-J model. Preferred-FT, on the other hand, shows minimal improvement over the original SFT model. A direct comparison using human raters, reported in Section~\ref{sec:human-judgments}, reveals that samples produced by DPO at a temperature of 0.25 are chosen over PPO's (at temperature 0.0) 58\% of the time.

For single-turn dialogue tasks, we assess the various approaches using the portion of the Anthropic HH dataset test split \citep{bai2022training} that involves one exchange between a human and an assistant. For the GPT-4 based assessments, we utilize the completions preferred by humans as references to calculate win rates across the different strategies. Given the absence of an established SFT model for this setting, our pipeline begins with a pre-trained Pythia-2.8B model. We apply Preferred-FT on the selected completions to fine-tune a reference model, ensuring generated completions remain within the training distribution, and subsequently apply DPO for further optimization. Our experimental setup also includes comparisons with the best of 128 Preferred-FT completions (with performance leveling off beyond $N = 128$ as detailed in Appendix Figure~\ref{fig:best-of-n}), as well as a 2-shot prompt adaptation of the original Pythia-2.8B model. Across all approaches, DPO either matches or surpasses others at their respective optimal sampling temperatures. Furthermore, we investigate an RLHF model trained with PPO on the Anthropic HH dataset \footnote{\url{https://huggingface.co/reciprocate/ppo_hh_pythia-6B}}, made publicly available \footnote{\url{https://github.com/CarperAI/trlx/tree/main/examples/hh}}, but are unable to achieve better results than those of the base Pythia-2.8B regardless of prompt selection or sampling temperature. In line with TL;DR outcomes and recognizing that both approaches target an equivalent reward objective, we interpret Best of 128 as an approximate indicator of PPO-level performance. Taken together, DPO stands out as the only method that is both computationally efficient and capable of outperforming the preferred human completions from the Anthropic HH dataset, while delivering results on par with or exceeding those of the computationally intensive Best of 128 method. Lastly, Figure~\ref{fig:dialogue-main} demonstrates that DPO attains peak performance after relatively few training steps.

\subsection{Generalization to a new input distribution}

To more thoroughly examine how PPO and DPO perform when faced with a distributional shift, we assess both policies—trained on Reddit TL;DR summarization—on out-of-distribution news articles from the CNN/DailyMail dataset's test split \citep{nallapati-etal-2016-abstractive}. We use the optimal sampling temperatures determined for TL;DR (0 and 0.25). The outcomes are displayed in Table~\ref{tab:ood}. We measured the win rate of GPT-4 compared to reference summaries from the dataset, employing the same GPT-4 (C) prompt as in the Reddit TL;DR setup but substituting “forum post” with “news article.” DPO consistently surpasses the PPO policy by a considerable margin on this novel dataset. These results suggest that, even without leveraging extra unlabeled Reddit TL;DR prompts as PPO does, DPO can generalize across distributions at least as effectively as PPO.

\subsection{Validating GPT-4 judgments with human judgments}
\label{sec:human-judgments}
We undertake a human evaluation to assess the accuracy of GPT-4’s judgments, drawing upon summaries generated in the TL;DR experiment and utilizing two distinct GPT-4 prompts. The \textbf{GPT-4 (S)} (simple) prompt requests the selection of the summary that best captures the salient information from the post. In contrast, the \textbf{GPT-4 (C)} (concise) prompt extends this by also soliciting a judgment on conciseness; this distinction is motivated by our observation that GPT-4 using the (S) prompt favors longer, more repetitive summaries than those preferred by human evaluators. The full text of the prompts is available in Appendix~\ref{app:prompts}. We perform three comparative evaluations—using DPO (temp. 0.25) as the best, PPO (temp. 1.0) as the worst, and SFT (temp. 0.25) as an intermediate performer—in each case comparing them against the PPO policy at its best sampling temperature (greedy sampling). These represent a range of sample quality. Our findings indicate that GPT-4’s assessments, for both prompts, are about as consistent with human ratings as the agreement observed among human raters themselves (for DPO and PPO-1, we gather multiple human ratings, constrained by rater availability). In sum, win rates derived from the \textbf{GPT-4 (C)} prompt most closely track human preferences, so we rely on this prompt for reporting key results in Section~\ref{sec:dpo-real-datasets}. Further details, including the rating interface and a roster of participating human volunteers, are documented in Appendix~\ref{app:human-study}.

\section{Discussion}
Preference-based learning offers an effective and scalable approach for developing capable and aligned language models. In this work, we have presented DPO, a straightforward method that enables training language models directly from preference data, circumventing the need for reinforcement learning. Instead of adapting the preference learning task to a typical RL framework to leverage standard RL algorithms, DPO establishes a direct correspondence between language model policies and reward functions. This correspondence permits the language model to be optimized according to human preferences using a cross-entropy loss function, thereby eliminating the necessity for reinforcement learning without sacrificing generality. Remarkably, DPO achieves comparable or superior results relative to leading RLHF algorithms, such as those utilizing PPO, even in the absence of extensive hyperparameter tuning. This substantially lowers the practical challenges of training language models using human preferences.

\textbf{Limitations \& Future Work.} Several avenues for future exploration emerge from our findings. One key question concerns how well the DPO-trained policy extrapolates to data outside the training distribution, particularly in comparison to models trained with explicit reward functions. Preliminary evidence indicates that DPO generalization capabilities are on par with models trained via PPO, yet a more thorough investigation is warranted. For instance, it remains to be studied whether DPO policies, when self-labeling, can leverage unlabeled prompts as effectively. Moreover, the dynamics of reward over-optimization in the context of direct preference optimization require further analysis, especially given the slight performance decline observed in Figure~\ref{fig:dialogue-main}-right and whether this is attributable to over-optimization. Furthermore, while our evaluation encompasses models with up to 6B parameters, expanding DPO to much larger, cutting-edge models constitutes a promising direction for future research. With respect to model assessment, our analysis reveals that the win rates assessed by GPT-4 are sensitive to prompt formulation; subsequent work may focus on optimizing prompt strategies to obtain more reliable automated evaluations. Lastly, the potential of DPO extends beyond human preference-based training for language models and invites application to the training of generative models across diverse domains.